% =========================================================
% Proof Stub: Residue-Class Convergence
% Source: volume-iii/analysis/sequences/notes/sequences/notes-subsequence-toolkit.tex
% Mode: standard
% =========================================================

\clearpage
\phantomsection
\hypertarget{proof-RA-ABB-C-residue-class-conv}{}

\subsubsection[Residue-Class Convergence]{Proof --- Residue-Class Convergence}

\bigskip

\noindent
\textbf{Source.}
See \texttt{notes-subsequence-toolkit.tex}.

\vspace{0.75em}

\noindent
\textbf{Goal.}
If for each $r \in \{0,\ldots,k-1\}$ the subsequence $(a_{kn+r})$ converges to $L$, then $(a_n) \to L$.

\vspace{0.75em}

\noindent
\textbf{Logical form.}
\[
  \forall r \in \{0,\ldots,k-1\},\; a_{kn+r} \to L \Rightarrow a_n \to L.
\]

\vspace{0.75em}

\noindent
\textbf{Proof strategy.}
The residue classes $\{kn+r : n \in \mathbb{N}\}$ for $r = 0,\ldots,k-1$ form a finite partition of $\mathbb{N}$. Apply the Finite Partition Convergence Principle.

\vspace{0.75em}

\noindent
\textbf{Proof.}
\begin{proof}
\Claim If for each $r \in \{0,\ldots,k-1\}$ the subsequence $(a_{kn+r})$ converges to $L$, then $(a_n) \to L$.

\Given 

\Goal 

\Strategy The residue classes $\{kn+r : n \in \mathbb{N}\}$ for $r = 0,\ldots,k-1$ form a finite partition of $\mathbb{N}$. Apply the Finite Partition Convergence Principle.

\medskip

% --- observe N = E_0 ∪ ... ∪ E_{k-1} where E_r = {kn+r : n ∈ N} ---
% --- this is a finite partition into infinite sets ---
% --- each (a_n)_{n ∈ E_r} = (a_{kn+r}) → L ---
% --- apply Finite Partition Convergence Principle ---

\end{proof}

\begin{remark*}[Proof shape]
Immediate from Finite Partition Convergence Principle applied to residue classes.
\end{remark*}

\begin{remark*}[Dependencies]
\begin{itemize}
  \item Finite Partition Convergence Principle.
  \item Residue classes partition $\mathbb{N}$.
\end{itemize}
\end{remark*}
